\documentclass[10pt,a4paper]{amsart}
\usepackage{amsmath,amssymb,amsthm,graphicx}
% Retained minimal preamble for the Stacks Project chapter more-algebra.tex (More on Algebra).
% The chapter file pinned-original-source.tex carries only \input{preamble} on its first line, so no
% preamble text lives inside the chapter itself.  The Stacks Project uses one single shared preamble for
% every chapter; the snapshot stages it identically as sources/stacks-chow/preamble.tex and
% sources/stacks-etale-cohomology/preamble.tex.  The block below is copied mechanically, line for line,
% from sources/stacks-chow/preamble.tex lines 14-18 (the Xy-pic setup) and 154-217 (the native theorem
% declarations together with the mathematical macros).  Nothing was added, removed or redefined; the
% wrapper class and the amsmath/amssymb/amsthm/graphicx base packages are supplied by the assembler.
% Chapter: more-algebra.tex -- More on Algebra.
% For commutative diagrams we use Xy-pic
\usepackage[all]{xy}

% We use 2cell for 2-commutative diagrams.
\xyoption{2cell}

% Theorem environments.
%
\theoremstyle{plain}
\newtheorem{theorem}[subsection]{Theorem}
\newtheorem{proposition}[subsection]{Proposition}
\newtheorem{lemma}[subsection]{Lemma}

\theoremstyle{definition}
\newtheorem{definition}[subsection]{Definition}
\newtheorem{example}[subsection]{Example}
\newtheorem{exercise}[subsection]{Exercise}
\newtheorem{situation}[subsection]{Situation}

\theoremstyle{remark}
\newtheorem{remark}[subsection]{Remark}
\newtheorem{remarks}[subsection]{Remarks}

\numberwithin{equation}{subsection}

% Macros
%
\def\lim{\mathop{\mathrm{lim}}\nolimits}
\def\colim{\mathop{\mathrm{colim}}\nolimits}
\def\Spec{\mathop{\mathrm{Spec}}}
\def\Hom{\mathop{\mathrm{Hom}}\nolimits}
\def\Ext{\mathop{\mathrm{Ext}}\nolimits}
\def\SheafHom{\mathop{\mathcal{H}\!\mathit{om}}\nolimits}
\def\SheafExt{\mathop{\mathcal{E}\!\mathit{xt}}\nolimits}
\def\Sch{\mathit{Sch}}
\def\Mor{\mathop{\mathrm{Mor}}\nolimits}
\def\Ob{\mathop{\mathrm{Ob}}\nolimits}
\def\Sh{\mathop{\mathit{Sh}}\nolimits}
\def\NL{\mathop{N\!L}\nolimits}
\def\CH{\mathop{\mathrm{CH}}\nolimits}
\def\proetale{{pro\text{-}\acute{e}tale}}
\def\etale{{\acute{e}tale}}
\def\QCoh{\mathit{QCoh}}
\def\Ker{\mathop{\mathrm{Ker}}}
\def\Im{\mathop{\mathrm{Im}}}
\def\Coker{\mathop{\mathrm{Coker}}}
\def\Coim{\mathop{\mathrm{Coim}}}

% Boxtimes
%
\DeclareMathSymbol{\boxtimes}{\mathbin}{AMSa}{"02}

%
% Macros for moduli stacks/spaces
%
\def\QCohstack{\mathcal{QC}\!\mathit{oh}}
\def\Cohstack{\mathcal{C}\!\mathit{oh}}
\def\Spacesstack{\mathcal{S}\!\mathit{paces}}
\def\Quotfunctor{\mathrm{Quot}}
\def\Hilbfunctor{\mathrm{Hilb}}
\def\Curvesstack{\mathcal{C}\!\mathit{urves}}
\def\Polarizedstack{\mathcal{P}\!\mathit{olarized}}
\def\Complexesstack{\mathcal{C}\!\mathit{omplexes}}
% \Pic is the operator that assigns to X its picard group, usage \Pic(X)
% \Picardstack_{X/B} denotes the Picard stack of X over B
% \Picardfunctor_{X/B} denotes the Picard functor of X over B
\def\Pic{\mathop{\mathrm{Pic}}\nolimits}
\def\Picardstack{\mathcal{P}\!\mathit{ic}}
\def\Picardfunctor{\mathrm{Pic}}
\def\Deformationcategory{\mathcal{D}\!\mathit{ef}}

% --- appended 2026-10-12 (Chief): the project own definitions of the reference/slogan/history
% display environments, copied mechanically from the shared Stacks preamble lines 8-12
% (\usepackage{verbatim} plus the three \newenvironment lines, which the chapter preamble keeps
% verbatim from the project preamble). Nothing added, removed or redefined: these environments
% carry no mathematical content and the project itself typesets them as comments. Without these
% three lines a counted statement whose own text uses them fails with Environment slogan undefined.
\usepackage{verbatim}
\newenvironment{reference}{\comment}{\endcomment}
\newenvironment{slogan}{\comment}{\endcomment}
\newenvironment{history}{\comment}{\endcomment}
\usepackage{cleveref}
\crefname{originallemma}{original lemma}{original lemmas}
\Crefname{originallemma}{Original lemma}{Original lemmas}
\makeatletter
\newlabel{more-algebra-lemma-henselization-NL}{{0D08}{source}{}{sourcepointer-1}{}}
\newlabel{more-algebra-lemma-henselization-NL@cref}{{[originallemma][0][]0D08}{[1][0][]source}{}{}{}}
\makeatother
\crefname{originallemma}{original lemma}{original lemmas}
\Crefname{originallemma}{Original lemma}{Original lemmas}
\makeatletter
\newlabel{algebra-lemma-henselian-functorial}{{04GS}{source}{}{sourcepointer-2}{}}
\newlabel{algebra-lemma-henselian-functorial@cref}{{[originallemma][0][]04GS}{[1][0][]source}{}{}{}}
\makeatother
\crefname{originallemma}{original lemma}{original lemmas}
\Crefname{originallemma}{Original lemma}{Original lemmas}
\makeatletter
\newlabel{algebra-lemma-strictly-henselian-functorial}{{04GU}{source}{}{sourcepointer-3}{}}
\newlabel{algebra-lemma-strictly-henselian-functorial@cref}{{[originallemma][0][]04GU}{[1][0][]source}{}{}{}}
\makeatother
\crefname{originallemma}{original lemma}{original lemmas}
\Crefname{originallemma}{Original lemma}{Original lemmas}
\makeatletter
\newlabel{algebra-lemma-colimits-NL}{{07BQ}{source}{}{sourcepointer-4}{}}
\newlabel{algebra-lemma-colimits-NL@cref}{{[originallemma][0][]07BQ}{[1][0][]source}{}{}{}}
\makeatother
\crefname{originallemma}{original lemma}{original lemmas}
\Crefname{originallemma}{Original lemma}{Original lemmas}
\makeatletter
\newlabel{algebra-definition-etale}{{00U1}{source}{}{sourcepointer-5}{}}
\newlabel{algebra-definition-etale@cref}{{[originallemma][0][]00U1}{[1][0][]source}{}{}{}}
\makeatother
\crefname{originallemma}{original lemma}{original lemmas}
\Crefname{originallemma}{Original lemma}{Original lemmas}
\makeatletter
\newlabel{algebra-lemma-exact-sequence-NL}{{00S2}{source}{}{sourcepointer-6}{}}
\newlabel{algebra-lemma-exact-sequence-NL@cref}{{[originallemma][0][]00S2}{[1][0][]source}{}{}{}}
\makeatother
\crefname{originallemma}{original lemma}{original lemmas}
\Crefname{originallemma}{Original lemma}{Original lemmas}
\makeatletter
\newlabel{algebra-lemma-etale-standard-smooth}{{00U9}{source}{}{sourcepointer-7}{}}
\newlabel{algebra-lemma-etale-standard-smooth@cref}{{[originallemma][0][]00U9}{[1][0][]source}{}{}{}}
\makeatother
\crefname{originallemma}{original lemma}{original lemmas}
\Crefname{originallemma}{Original lemma}{Original lemmas}
\makeatletter
\newlabel{lemma-transitive-colimit-lci-at-end}{{07D5}{source}{}{sourcepointer-8}{}}
\newlabel{lemma-transitive-colimit-lci-at-end@cref}{{[originallemma][0][]07D5}{[1][0][]source}{}{}{}}
\makeatother
\title{The naive cotangent complex of henselization and strict henselization}
\author{The Stacks Project Authors}
\date{stacks-project:more-algebra.tex:local-snapshot-20261012; GNU FDL 1.2 or later}
\begin{document}
\maketitle
\setcounter{subsection}{1}
\noindent\textit{Editorial note.} All delivered mathematical text is the author's own, copied from the pinned Stacks Project chapter file more-algebra.tex; no statement or proof step was rewritten, reordered or removed.\par\medskip
\noindent\textit{Editorial note.} The chapter file pinned-original-source.tex carries only \textbackslash{}input\{preamble\} on its first line, so no preamble text lives inside the chapter. The retained file sources/stacks-more-algebra/retained-preamble.tex is a mechanical, line-for-line copy of the project's single shared preamble as staged in the snapshot (Xy-pic setup plus theorem declarations and mathematical macros). Nothing was added, removed or redefined, and the excerpt serves only as preamble, never as a span coordinate.\par\medskip
\noindent\textit{Editorial note.} Both counted spans start at the \textbackslash{}begin line of their environment, so the statement and the author's proof are delivered whole; surrounding prose of the section is deliberately outside the counted spans and is carried as standard prerequisites instead.\par\medskip
\noindent\textit{Editorial note.} The citation fields (source\_id, source\_version, authors, source\_url, license.name, source\_title) are identical to those of every other card this lane has produced from this chapter, so the batch selector treats them as one source.\par\medskip
\noindent\textit{Allowed standard prerequisites (editorial).}\begin{itemize}
\item The relative naive cotangent complex of a filtered colimit of etale algebras is acyclic. Algebra, Lemma \textbackslash{}ref\{algebra-lemma-colimits-NL\} and Algebra, Definition \textbackslash{}ref\{algebra-definition-etale\}, cited at proof lines 8536-8539. Makes the intermediate naive cotangent complexes vanish on cohomology.
\item The Jacobi-Zariski sequence for a chain of ring maps. Algebra, Lemma \textbackslash{}ref\{algebra-lemma-exact-sequence-NL\}, cited at proof lines 8540-8542. Removes the henselization of the source from the comparison.
\item An etale ring map is a global complete intersection, and the six-term Jacobi-Zariski sequence for a colimit of local complete intersections. Algebra, Lemma \textbackslash{}ref\{algebra-lemma-etale-standard-smooth\} and Lemma \textbackslash{}ref\{lemma-transitive-colimit-lci-at-end\}, cited at proof lines 8545-8548. Produces the exact sequence that identifies the base change of the naive cotangent complex with the one over the henselization.
\end{itemize}\medskip
\ifdefined\hypertarget\hypertarget{sourcepointer-1}{}\fi
\noindent\textit{Source reference.} Original lemma 0D08: focused-seeds/AGP/sources/stacks-snapshot-20261011/tags.txt, source lines 16303--16303. The lemma's own canonical tag, quoted so that readers can resolve the label inside the delivered text; it points at the counted statement itself and carries no proof content.\par\medskip
\ifdefined\hypertarget\hypertarget{sourcepointer-2}{}\fi
\noindent\textit{Source reference.} Original lemma 04GS: focused-seeds/AGP/sources/stacks-snapshot-20261011/tags.txt, source lines 5843--5843. Navigation pointer to a statement the counted proof refers to but does not contain; nothing load-bearing is delegated to it.\par\medskip
\ifdefined\hypertarget\hypertarget{sourcepointer-3}{}\fi
\noindent\textit{Source reference.} Original lemma 04GU: focused-seeds/AGP/sources/stacks-snapshot-20261011/tags.txt, source lines 5845--5845. Navigation pointer to a statement the counted proof refers to but does not contain; nothing load-bearing is delegated to it.\par\medskip
\ifdefined\hypertarget\hypertarget{sourcepointer-4}{}\fi
\noindent\textit{Source reference.} Original lemma 07BQ: focused-seeds/AGP/sources/stacks-snapshot-20261011/tags.txt, source lines 9350--9350. Navigation pointer to a statement the counted proof refers to but does not contain; nothing load-bearing is delegated to it.\par\medskip
\ifdefined\hypertarget\hypertarget{sourcepointer-5}{}\fi
\noindent\textit{Source reference.} Original lemma 00U1: focused-seeds/AGP/sources/stacks-snapshot-20261011/tags.txt, source lines 1095--1095. Navigation pointer to a statement the counted proof refers to but does not contain; nothing load-bearing is delegated to it.\par\medskip
\ifdefined\hypertarget\hypertarget{sourcepointer-6}{}\fi
\noindent\textit{Source reference.} Original lemma 00S2: focused-seeds/AGP/sources/stacks-snapshot-20261011/tags.txt, source lines 1024--1024. Navigation pointer to a statement the counted proof refers to but does not contain; nothing load-bearing is delegated to it.\par\medskip
\ifdefined\hypertarget\hypertarget{sourcepointer-7}{}\fi
\noindent\textit{Source reference.} Original lemma 00U9: focused-seeds/AGP/sources/stacks-snapshot-20261011/tags.txt, source lines 1103--1103. Navigation pointer to a statement the counted proof refers to but does not contain; nothing load-bearing is delegated to it.\par\medskip
\ifdefined\hypertarget\hypertarget{sourcepointer-8}{}\fi
\noindent\textit{Source reference.} Original lemma 07D5: focused-seeds/AGP/sources/stacks-snapshot-20261011/tags.txt, source lines 9400--9400. Navigation pointer to a statement the counted proof refers to but does not contain; nothing load-bearing is delegated to it.\par\medskip
\begin{lemma}
\label{lemma-henselization-NL}
Let $A \to B$ be a local homomorphism of local rings.
Let $A^h \to B^h$, resp.\ $A^{sh} \to B^{sh}$ be the induced
map on henselizations, resp.\ strict henselizations
(Algebra, Lemma \ref{algebra-lemma-henselian-functorial},
resp.\ Lemma \ref{algebra-lemma-strictly-henselian-functorial}).
Then $\NL_{B/A} \otimes_B B^h \to \NL_{B^h/A^h}$ and
$\NL_{B/A} \otimes_B B^{sh} \to \NL_{B^{sh}/A^{sh}}$
induce isomorphisms on cohomology groups.
\end{lemma}

\begin{proof}
Since $A^h$ is a filtered colimit of \'etale algebras over $A$
we see that $\NL_{A^h/A}$ is an acyclic complex by
Algebra, Lemma \ref{algebra-lemma-colimits-NL} and
Algebra, Definition \ref{algebra-definition-etale}.
The same is true for $B^h/B$.
Using the Jacobi-Zariski sequence
(Algebra, Lemma \ref{algebra-lemma-exact-sequence-NL})
for $A \to A^h \to B^h$ we find that
$\NL_{B^h/A} \to \NL_{B^h/A^h}$ induces isomorphisms
on cohomology groups. Moreover, an \'etale ring map is
a local complete intersection as it is even a global complete
intersection, see Algebra, Lemma \ref{algebra-lemma-etale-standard-smooth}.
By Lemma \ref{lemma-transitive-colimit-lci-at-end}
we get a six term exact Jacobi-Zariski sequence
associated to $A \to B \to B^h$ which proves that
$\NL_{B/A} \otimes_B B^h \to \NL_{B^h/A}$ induces isomorphisms
on cohomology groups. This finishes the proof in the
case of the map on henselizations. The case of strict henselization
is proved in exactly the same manner.
\end{proof}

\end{document}
